% --- Archon physics formalization source begin ---
% archon:physics
% archon:covers PhyXMiniProblems/problem_phyx_mini_0040.lean
% archon:source-report reports/phyx_mini/problem_phyx_mini_0040.source.json

\chapter{Physics problem phyx\_mini\_0040}
\label{ch:PhyXMiniProblems_problem_phyx_mini_0040}

\paragraph{Problem source.}
\#\# Physical scenario

A cylindrical glass rod (figure) has index of refraction 1.52. It is surrounded by air. One end is ground to a hemispherical surface with radius \textbackslash{}( R = 2.00 \textbackslash{} \textbackslash{}mathrm\{cm\} \textbackslash{}). A small object is placed on the axis of the rod, 8.00 \textbackslash{} \textbackslash{}mathrm\{cm\} to the left of the vertex.

\#\# Question

Find the lateral magnification.

\#\# Answer choices

- A: A: -0.856

- B: B: -0.929

- C: C: +0.995

- D: D: -0.814

\#\# Figure caption (auxiliary; use the image as primary evidence)

The image illustrates a refraction scenario involving light passing through a lens or curved surface.

- **Medium and Indices of Refraction:**
  - The left region is marked with \textbackslash{}( n\_a = 1.00 \textbackslash{}) labeled as air.
  - The right region is marked with \textbackslash{}( n\_b = 1.52 \textbackslash{}), indicating a different medium.

- **Points:**
  - Point \textbackslash{}( P \textbackslash{}): Represented by a blue dot on the left.
  - Point \textbackslash{}( P' \textbackslash{}): Represented by a red dot on the right.
  - Point \textbackslash{}( C \textbackslash{}): Represented by a black dot at the center of the curved surface.

- **Rays:**
  - Three purple rays originate from point \textbackslash{}( P \textbackslash{}) and converge towards the curved surface, then continue to point \textbackslash{}( P' \textbackslash{}).

- **Curved Surface:**
  - Curved surface is depicted with a bulge to the right.

- **Measurements:**
  - Distance from \textbackslash{}( P \textbackslash{}) to the curved surface \textbackslash{}( s = 8.00 \textbackslash{}, \textbackslash{}text\{cm\} \textbackslash{}).
  - Distance from the curved surface to \textbackslash{}( P' \textbackslash{}) \textbackslash{}( s' \textbackslash{}).
  - Radius of curvature \textbackslash{}( R = 2.00 \textbackslash{}, \textbackslash{}text\{cm\} \textbackslash{}).

This setup is typically used to demonstrate the refraction of light between two different media.

\#\# Dataset metadata

Geometrical Optics; Physical Model Grounding Reasoning, Multi-Formula Reasoning

\paragraph{Recorded answer/context.}
B: B: -0.929

\paragraph{Figure/image path.}
/root/proposal\_for\_physic/science-mango/phyx\_mini\_run/phyx\_data/test\_image/40.png

\paragraph{Formalization target.}
create a compiling Lean file with sorry bodies at `PhyXMiniProblems/problem\_phyx\_mini\_0040.lean`. The Lean declarations must preserve the physical quantities, dimensions or dimensional roles, figure labels, governing-law hypotheses, and final relation expressed by this problem.
Use Mathlib/Physlib names found through LeanExplore where available. If a physics API is missing, introduce faithful local abstractions rather than scalar placeholder aliases.

\begin{theorem}[Physics formalization target]
\label{thm:physics:phyx_mini_0040:target}
\lean{PhyXMiniProblems.ProblemPhyXMini0040.problem_phyx_mini_0040}
\uses{def:physics:phyx-mini-0040:phyxminiproblems-problemphyxmini0040-hemisphericalrodsetup, def:physics:phyx-mini-0040:phyxminiproblems-problemphyxmini0040-matchesfigurereadout, def:physics:phyx-mini-0040:phyxminiproblems-problemphyxmini0040-hasphysicalfiguregeometry, def:physics:phyx-mini-0040:phyxminiproblems-problemphyxmini0040-hasfirstorderparaxialimage, def:physics:phyx-mini-0040:phyxminiproblems-problemphyxmini0040-transverseimagemap, def:physics:phyx-mini-0040:phyxminiproblems-problemphyxmini0040-hasfirstorderlateralmagnification, def:physics:phyx-mini-0040:phyxminiproblems-problemphyxmini0040-isuniqueclosestanswer}
The assigned autoformalize agent should translate this physics problem into Lean declarations in the covered file, with theorem and lemma proof bodies written as `by sorry`.
\end{theorem}
\begin{proof}
This is an autoformalization task, not a proof task. Produce statements that can later be proved by the physics prover without weakening the physical model.
\end{proof}

% --- Archon declaration topology begin ---
\section{Declaration topology}
\begin{definition}[Optical Length]
  \label{def:physics:phyx-mini-0040:phyxminiproblems-problemphyxmini0040-opticallength}
  \lean{PhyXMiniProblems.ProblemPhyXMini0040.OpticalLength}
  A signed physical length whose readout changes coherently with unit choice
\end{definition}
\begin{proof}
  This is the declared model definition of Optical Length.
\end{proof}
\begin{definition}[length In]
  \label{def:physics:phyx-mini-0040:phyxminiproblems-problemphyxmini0040-lengthin}
  \lean{PhyXMiniProblems.ProblemPhyXMini0040.lengthIn}
  \uses{def:physics:phyx-mini-0040:phyxminiproblems-problemphyxmini0040-opticallength}
  The physical length having scalar readout magnitude in the chosen unit
\end{definition}
\begin{proof}
  This is the declared model definition of length In.
\end{proof}
\begin{definition}[length Readout]
  \label{def:physics:phyx-mini-0040:phyxminiproblems-problemphyxmini0040-lengthreadout}
  \lean{PhyXMiniProblems.ProblemPhyXMini0040.lengthReadout}
  \uses{def:physics:phyx-mini-0040:phyxminiproblems-problemphyxmini0040-opticallength}
  The scalar readout of a physical length in the selected length unit
\end{definition}
\begin{proof}
  This is the declared model definition of length Readout.
\end{proof}
\begin{definition}[Optical Medium]
  \label{def:physics:phyx-mini-0040:phyxminiproblems-problemphyxmini0040-opticalmedium}
  \lean{PhyXMiniProblems.ProblemPhyXMini0040.OpticalMedium}
  The two homogeneous optical media on the incident and transmitted sides
\end{definition}
\begin{proof}
  This is the declared model definition of Optical Medium.
\end{proof}
\begin{definition}[Axial Point]
  \label{def:physics:phyx-mini-0040:phyxminiproblems-problemphyxmini0040-axialpoint}
  \lean{PhyXMiniProblems.ProblemPhyXMini0040.AxialPoint}
  Axial point labels appearing in the source figure
\end{definition}
\begin{proof}
  This is the declared model definition of Axial Point.
\end{proof}
\begin{definition}[Rod Shape]
  \label{def:physics:phyx-mini-0040:phyxminiproblems-problemphyxmini0040-rodshape}
  \lean{PhyXMiniProblems.ProblemPhyXMini0040.RodShape}
  The overall shape of the transparent solid in the physical scenario
\end{definition}
\begin{proof}
  This is the declared model definition of Rod Shape.
\end{proof}
\begin{definition}[End Surface Shape]
  \label{def:physics:phyx-mini-0040:phyxminiproblems-problemphyxmini0040-endsurfaceshape}
  \lean{PhyXMiniProblems.ProblemPhyXMini0040.EndSurfaceShape}
  The shape ground onto the incident end of the glass rod
\end{definition}
\begin{proof}
  This is the declared model definition of End Surface Shape.
\end{proof}
\begin{definition}[Hemispherical Rod Setup]
  \label{def:physics:phyx-mini-0040:phyxminiproblems-problemphyxmini0040-hemisphericalrodsetup}
  \lean{PhyXMiniProblems.ProblemPhyXMini0040.HemisphericalRodSetup}
  \uses{def:physics:phyx-mini-0040:phyxminiproblems-problemphyxmini0040-opticallength, def:physics:phyx-mini-0040:phyxminiproblems-problemphyxmini0040-opticalmedium, def:physics:phyx-mini-0040:phyxminiproblems-problemphyxmini0040-axialpoint, def:physics:phyx-mini-0040:phyxminiproblems-problemphyxmini0040-rodshape, def:physics:phyx-mini-0040:phyxminiproblems-problemphyxmini0040-endsurfaceshape}
  Physical and geometric quantities attached to the air--glass refraction diagram. The axial coordinate is signed, while surfaceRadius is a positive physical length when the setup is physically admissible
\end{definition}
\begin{proof}
  This is the declared model definition of Hemispherical Rod Setup.
\end{proof}
\begin{definition}[axial Separation Readout]
  \label{def:physics:phyx-mini-0040:phyxminiproblems-problemphyxmini0040-axialseparationreadout}
  \lean{PhyXMiniProblems.ProblemPhyXMini0040.axialSeparationReadout}
  \uses{def:physics:phyx-mini-0040:phyxminiproblems-problemphyxmini0040-lengthreadout, def:physics:phyx-mini-0040:phyxminiproblems-problemphyxmini0040-axialpoint, def:physics:phyx-mini-0040:phyxminiproblems-problemphyxmini0040-hemisphericalrodsetup}
  Signed axial separation from start to finish, read in the selected unit
\end{definition}
\begin{proof}
  This is the declared model definition of axial Separation Readout.
\end{proof}
\begin{definition}[object Distance Readout]
  \label{def:physics:phyx-mini-0040:phyxminiproblems-problemphyxmini0040-objectdistancereadout}
  \lean{PhyXMiniProblems.ProblemPhyXMini0040.objectDistanceReadout}
  \uses{def:physics:phyx-mini-0040:phyxminiproblems-problemphyxmini0040-hemisphericalrodsetup, def:physics:phyx-mini-0040:phyxminiproblems-problemphyxmini0040-axialseparationreadout}
  Figure label s: positive object distance from P to the surface vertex
\end{definition}
\begin{proof}
  This is the declared model definition of object Distance Readout.
\end{proof}
\begin{definition}[image Distance Readout]
  \label{def:physics:phyx-mini-0040:phyxminiproblems-problemphyxmini0040-imagedistancereadout}
  \lean{PhyXMiniProblems.ProblemPhyXMini0040.imageDistanceReadout}
  \uses{def:physics:phyx-mini-0040:phyxminiproblems-problemphyxmini0040-hemisphericalrodsetup, def:physics:phyx-mini-0040:phyxminiproblems-problemphyxmini0040-axialseparationreadout}
  Figure label s': positive real-image distance from the vertex to P'
\end{definition}
\begin{proof}
  This is the declared model definition of image Distance Readout.
\end{proof}
\begin{definition}[radius Readout]
  \label{def:physics:phyx-mini-0040:phyxminiproblems-problemphyxmini0040-radiusreadout}
  \lean{PhyXMiniProblems.ProblemPhyXMini0040.radiusReadout}
  \uses{def:physics:phyx-mini-0040:phyxminiproblems-problemphyxmini0040-lengthreadout, def:physics:phyx-mini-0040:phyxminiproblems-problemphyxmini0040-hemisphericalrodsetup}
  Scalar readout of the hemispherical surface's physical radius
\end{definition}
\begin{proof}
  This is the declared model definition of radius Readout.
\end{proof}
\begin{definition}[Matches Figure Readout]
  \label{def:physics:phyx-mini-0040:phyxminiproblems-problemphyxmini0040-matchesfigurereadout}
  \lean{PhyXMiniProblems.ProblemPhyXMini0040.MatchesFigureReadout}
  \uses{def:physics:phyx-mini-0040:phyxminiproblems-problemphyxmini0040-hemisphericalrodsetup, def:physics:phyx-mini-0040:phyxminiproblems-problemphyxmini0040-axialseparationreadout, def:physics:phyx-mini-0040:phyxminiproblems-problemphyxmini0040-objectdistancereadout, def:physics:phyx-mini-0040:phyxminiproblems-problemphyxmini0040-radiusreadout}
  Problem and figure readouts. They fix the rod and end shapes, the two dimensionless refractive indices, s = 8.00 cm, and R = 2.00 cm. The relation between the vertex and C records that C is the center of the hemisphere. No numerical value is assigned to the derived image distance
\end{definition}
\begin{proof}
  This is the declared model definition of Matches Figure Readout.
\end{proof}
\begin{definition}[Has Physical Figure Geometry]
  \label{def:physics:phyx-mini-0040:phyxminiproblems-problemphyxmini0040-hasphysicalfiguregeometry}
  \lean{PhyXMiniProblems.ProblemPhyXMini0040.HasPhysicalFigureGeometry}
  \uses{def:physics:phyx-mini-0040:phyxminiproblems-problemphyxmini0040-lengthreadout, def:physics:phyx-mini-0040:phyxminiproblems-problemphyxmini0040-hemisphericalrodsetup, def:physics:phyx-mini-0040:phyxminiproblems-problemphyxmini0040-radiusreadout}
  Physical branch information visible in the diagram: refractive indices and radius are positive, and the labels occur in the order P, vertex, C, P' along the propagation axis
\end{definition}
\begin{proof}
  This is the declared model definition of Has Physical Figure Geometry.
\end{proof}
\begin{definition}[surface Axial Offset Readout]
  \label{def:physics:phyx-mini-0040:phyxminiproblems-problemphyxmini0040-surfaceaxialoffsetreadout}
  \lean{PhyXMiniProblems.ProblemPhyXMini0040.surfaceAxialOffsetReadout}
  \uses{def:physics:phyx-mini-0040:phyxminiproblems-problemphyxmini0040-hemisphericalrodsetup, def:physics:phyx-mini-0040:phyxminiproblems-problemphyxmini0040-radiusreadout}
  At transverse height rayHeight, this is the axial offset of the incident point on the left-hand hemispherical surface. The formula is the exact circle geometry R - sqrt (R\textasciicircum{}2 - h\textasciicircum{}2), not a paraxial replacement
\end{definition}
\begin{proof}
  This is the declared model definition of surface Axial Offset Readout.
\end{proof}
\begin{definition}[incident Ray Angle]
  \label{def:physics:phyx-mini-0040:phyxminiproblems-problemphyxmini0040-incidentrayangle}
  \lean{PhyXMiniProblems.ProblemPhyXMini0040.incidentRayAngle}
  \uses{def:physics:phyx-mini-0040:phyxminiproblems-problemphyxmini0040-hemisphericalrodsetup, def:physics:phyx-mini-0040:phyxminiproblems-problemphyxmini0040-objectdistancereadout, def:physics:phyx-mini-0040:phyxminiproblems-problemphyxmini0040-surfaceaxialoffsetreadout}
  Exact direction angle of the ray from the axial object P to the surface
\end{definition}
\begin{proof}
  This is the declared model definition of incident Ray Angle.
\end{proof}
\begin{definition}[inward Normal Angle]
  \label{def:physics:phyx-mini-0040:phyxminiproblems-problemphyxmini0040-inwardnormalangle}
  \lean{PhyXMiniProblems.ProblemPhyXMini0040.inwardNormalAngle}
  \uses{def:physics:phyx-mini-0040:phyxminiproblems-problemphyxmini0040-hemisphericalrodsetup, def:physics:phyx-mini-0040:phyxminiproblems-problemphyxmini0040-radiusreadout}
  Exact direction angle of the inward surface normal. At a point above the axis it points downwards from the spherical surface toward C
\end{definition}
\begin{proof}
  This is the declared model definition of inward Normal Angle.
\end{proof}
\begin{definition}[image Ray Angle]
  \label{def:physics:phyx-mini-0040:phyxminiproblems-problemphyxmini0040-imagerayangle}
  \lean{PhyXMiniProblems.ProblemPhyXMini0040.imageRayAngle}
  \uses{def:physics:phyx-mini-0040:phyxminiproblems-problemphyxmini0040-hemisphericalrodsetup, def:physics:phyx-mini-0040:phyxminiproblems-problemphyxmini0040-imagedistancereadout, def:physics:phyx-mini-0040:phyxminiproblems-problemphyxmini0040-surfaceaxialoffsetreadout}
  Exact direction angle from the surface point to the proposed axial image P'
\end{definition}
\begin{proof}
  This is the declared model definition of image Ray Angle.
\end{proof}
\begin{definition}[axial Snell Residual]
  \label{def:physics:phyx-mini-0040:phyxminiproblems-problemphyxmini0040-axialsnellresidual}
  \lean{PhyXMiniProblems.ProblemPhyXMini0040.axialSnellResidual}
  \uses{def:physics:phyx-mini-0040:phyxminiproblems-problemphyxmini0040-hemisphericalrodsetup, def:physics:phyx-mini-0040:phyxminiproblems-problemphyxmini0040-incidentrayangle, def:physics:phyx-mini-0040:phyxminiproblems-problemphyxmini0040-inwardnormalangle, def:physics:phyx-mini-0040:phyxminiproblems-problemphyxmini0040-imagerayangle}
  Difference between the two sides of exact Snell's law for the incident ray and the ray aimed at P'. All angles are signed radian readouts relative to the propagation axis
\end{definition}
\begin{proof}
  This is the declared model definition of axial Snell Residual.
\end{proof}
\begin{definition}[Has First Order Paraxial Image]
  \label{def:physics:phyx-mini-0040:phyxminiproblems-problemphyxmini0040-hasfirstorderparaxialimage}
  \lean{PhyXMiniProblems.ProblemPhyXMini0040.HasFirstOrderParaxialImage}
  \uses{def:physics:phyx-mini-0040:phyxminiproblems-problemphyxmini0040-hemisphericalrodsetup, def:physics:phyx-mini-0040:phyxminiproblems-problemphyxmini0040-axialsnellresidual}
  P' is the first-order, or paraxial, image selected by exact spherical geometry and Snell's law. The derivative condition says that the exact Snell residual is o(rayHeight) at the optical axis. It does not assert exact stigmatic imaging for finite-height rays, which a spherical surface generally does not provide
\end{definition}
\begin{proof}
  This is the declared model definition of Has First Order Paraxial Image.
\end{proof}
\begin{definition}[Transverse Image Map]
  \label{def:physics:phyx-mini-0040:phyxminiproblems-problemphyxmini0040-transverseimagemap}
  \lean{PhyXMiniProblems.ProblemPhyXMini0040.TransverseImageMap}
  \uses{def:physics:phyx-mini-0040:phyxminiproblems-problemphyxmini0040-opticallength}
  A physical map from transverse object displacement to image displacement
\end{definition}
\begin{proof}
  This is the declared model definition of Transverse Image Map.
\end{proof}
\begin{definition}[transverse Image Readout]
  \label{def:physics:phyx-mini-0040:phyxminiproblems-problemphyxmini0040-transverseimagereadout}
  \lean{PhyXMiniProblems.ProblemPhyXMini0040.transverseImageReadout}
  \uses{def:physics:phyx-mini-0040:phyxminiproblems-problemphyxmini0040-lengthin, def:physics:phyx-mini-0040:phyxminiproblems-problemphyxmini0040-lengthreadout, def:physics:phyx-mini-0040:phyxminiproblems-problemphyxmini0040-transverseimagemap}
  Scalar readout, in one length unit, of a physical transverse image map. Its derivative is dimensionless because input and output have the same physical dimension
\end{definition}
\begin{proof}
  This is the declared model definition of transverse Image Readout.
\end{proof}
\begin{definition}[vertex Snell Residual]
  \label{def:physics:phyx-mini-0040:phyxminiproblems-problemphyxmini0040-vertexsnellresidual}
  \lean{PhyXMiniProblems.ProblemPhyXMini0040.vertexSnellResidual}
  \uses{def:physics:phyx-mini-0040:phyxminiproblems-problemphyxmini0040-hemisphericalrodsetup, def:physics:phyx-mini-0040:phyxminiproblems-problemphyxmini0040-objectdistancereadout, def:physics:phyx-mini-0040:phyxminiproblems-problemphyxmini0040-imagedistancereadout, def:physics:phyx-mini-0040:phyxminiproblems-problemphyxmini0040-transverseimagemap, def:physics:phyx-mini-0040:phyxminiproblems-problemphyxmini0040-transverseimagereadout}
  Exact Snell residual for the chief ray that passes through the surface vertex. The incident ray starts at the displaced object point, and the transmitted ray ends at the displacement supplied by imageMap in the plane through P'
\end{definition}
\begin{proof}
  This is the declared model definition of vertex Snell Residual.
\end{proof}
\begin{definition}[Has First Order Lateral Magnification]
  \label{def:physics:phyx-mini-0040:phyxminiproblems-problemphyxmini0040-hasfirstorderlateralmagnification}
  \lean{PhyXMiniProblems.ProblemPhyXMini0040.HasFirstOrderLateralMagnification}
  \uses{def:physics:phyx-mini-0040:phyxminiproblems-problemphyxmini0040-hemisphericalrodsetup, def:physics:phyx-mini-0040:phyxminiproblems-problemphyxmini0040-transverseimagemap, def:physics:phyx-mini-0040:phyxminiproblems-problemphyxmini0040-transverseimagereadout, def:physics:phyx-mini-0040:phyxminiproblems-problemphyxmini0040-vertexsnellresidual}
  The local/asymptotic contract defining signed lateral magnification. The map fixes the optical axis, has derivative lateralMagnification there in every unit system, and satisfies exact Snell refraction through the vertex to first order. Consequently the familiar formula m = -nₐ s' / (n\_b s) is a theorem about the derivative, rather than a globally exact finite-angle premise. A negative derivative represents image inversion
\end{definition}
\begin{proof}
  This is the declared model definition of Has First Order Lateral Magnification.
\end{proof}
\begin{definition}[Answer Choice]
  \label{def:physics:phyx-mini-0040:phyxminiproblems-problemphyxmini0040-answerchoice}
  \lean{PhyXMiniProblems.ProblemPhyXMini0040.AnswerChoice}
  Labels of the four displayed multiple-choice answers
\end{definition}
\begin{proof}
  This is the declared model definition of Answer Choice.
\end{proof}
\begin{definition}[answer Magnification]
  \label{def:physics:phyx-mini-0040:phyxminiproblems-problemphyxmini0040-answermagnification}
  \lean{PhyXMiniProblems.ProblemPhyXMini0040.answerMagnification}
  \uses{def:physics:phyx-mini-0040:phyxminiproblems-problemphyxmini0040-answerchoice}
  Dimensionless lateral-magnification value printed beside each choice
\end{definition}
\begin{proof}
  This is the declared model definition of answer Magnification.
\end{proof}
\begin{definition}[Is Unique Closest Answer]
  \label{def:physics:phyx-mini-0040:phyxminiproblems-problemphyxmini0040-isuniqueclosestanswer}
  \lean{PhyXMiniProblems.ProblemPhyXMini0040.IsUniqueClosestAnswer}
  \uses{def:physics:phyx-mini-0040:phyxminiproblems-problemphyxmini0040-answerchoice, def:physics:phyx-mini-0040:phyxminiproblems-problemphyxmini0040-answermagnification}
  A choice is uniquely closest to the exact physical value displayed by the model
\end{definition}
\begin{proof}
  This is the declared model definition of Is Unique Closest Answer.
\end{proof}
\begin{lemma}[image Distance In Centimeters eq]
  \label{lem:physics:phyx-mini-0040:phyxminiproblems-problemphyxmini0040-imagedistanceincentimeters-eq}
  \lean{PhyXMiniProblems.ProblemPhyXMini0040.imageDistanceInCentimeters_eq}
  \uses{def:physics:phyx-mini-0040:phyxminiproblems-problemphyxmini0040-hemisphericalrodsetup, def:physics:phyx-mini-0040:phyxminiproblems-problemphyxmini0040-imagedistancereadout, def:physics:phyx-mini-0040:phyxminiproblems-problemphyxmini0040-matchesfigurereadout, def:physics:phyx-mini-0040:phyxminiproblems-problemphyxmini0040-hasphysicalfiguregeometry, def:physics:phyx-mini-0040:phyxminiproblems-problemphyxmini0040-hasfirstorderparaxialimage}
  The spherical-interface equation determines the derived image distance s' = 304/27 cm; this value is not included among the problem data
\end{lemma}
\begin{proof}
  The conclusion follows from the governing relations and figure readouts named in the statement of image Distance In Centimeters eq.
\end{proof}
% --- Archon declaration topology end ---
% --- Archon physics formalization source end ---
